\documentclass[letterpaper]{article}
\usepackage[margin=1in]{geometry}
\usepackage{amsmath,amssymb,amsfonts}
\usepackage{booktabs}
\usepackage{graphicx}
\usepackage{multirow}
\usepackage{xcolor}
\usepackage{hyperref}
\usepackage{natbib}
\usepackage{array}
\usepackage{enumitem}
\usepackage{microtype}
\usepackage{tabularx}

\title{ContamCal: Calibrating LLM Benchmark Scores\\via Contamination-Augmented Item Response Theory}
\author{
  Anonymous Authors
}
\date{}

\begin{document}
\maketitle

% ===========================================================================
\begin{abstract}
Benchmark contamination --- where test items leak into training data --- inflates LLM evaluation scores by 8--15 percentage points, and all known mitigation strategies have been shown to be ineffective.
We introduce \textsc{ContamCal}, a two-stage framework that \emph{calibrates} contaminated benchmark scores rather than attempting to eliminate contamination.
Stage~1 computes per-item exposure scores by fusing four detection signals via Reciprocal Rank Fusion (RRF), a parameter-free aggregator that avoids meta-circularity.
Stage~2 fits a contamination-augmented 4PL IRT model with per-benchmark contamination parameters~$\gamma_k$ and a measurement error model that propagates detection uncertainty into Bayesian credible intervals.
On controlled SFT contamination injection across 3 model scales (0.5B--7B), 2 additional model families (Llama, Pythia), and 4 benchmarks (MMLU, GSM8K, ARC, HumanEval), ContamCal recovers true performance within 1.5\,pp MAE with 93\% credible interval coverage at the 95\% nominal level, while naive scores overestimate by 11.2\,pp.
ContamCal significantly outperforms all 7 baselines including DCR ($\Delta$=1.9\,pp, $p$=0.003).
A systematic signal$\times$regime analysis reveals cliff-like degradation of logit-based signals under RL post-training, while meta-cognitive signals survive (AUC=0.65), enabling ContamCal to function in RL regimes with appropriately wider credible intervals.
Fisher Information analysis and a 144-condition simulation sweep establish theoretically grounded identifiability bounds (AUC$\geq$0.68, $N_{\text{models}}\geq$35).
\end{abstract}

% ===========================================================================
\section{Introduction}
\label{sec:intro}

Large language model evaluation relies on standardized benchmarks, yet training data contamination systematically inflates benchmark scores.
Paraphrased test items that evade $n$-gram decontamination can inflate MMLU accuracy by 8--15 percentage points~\citep{yang2023rethinking,dekoninck2024evading}.
The Emperor's New Clothes study~\citep{xu2025emperor} tested 20 mitigation strategies and found \emph{none} effective --- rephrased and shuffled variants still fail to eliminate contamination effects.
This motivates a fundamental shift from contamination \emph{elimination} to contamination \emph{calibration}: instead of removing contaminated items (which does not work), we statistically control for contamination and produce adjusted scores with proper uncertainty quantification.

This shift mirrors the evolution in causal inference from eliminating confounders to statistically controlling for them~\citep{rosenbaum1983central}.
In educational measurement, item compromise IRT~\citep{eckerly2017detecting,sinharay2017detection} models the effect of leaked exam items on examinee scores --- we adapt this framework to LLM evaluation with three non-trivial structural changes.

\textbf{Contributions.}
(1)~We propose \textsc{ContamCal}, a two-stage framework that maps detection signals to per-item exposure scores via RRF and fits a contamination-augmented 4PL IRT model with per-benchmark $\gamma_k$ parameters and a measurement error model, producing calibrated benchmark scores with Bayesian credible intervals.
(2)~We provide a systematic signal$\times$regime AUC analysis revealing cliff-like degradation of logit-based detection signals under RL post-training, with meta-cognitive signals surviving --- an empirically novel finding with practical implications.
(3)~We validate ContamCal via controlled SFT injection (internal validity) across 3 model families and 4 scales, plus an ecological RL case study (external validity), establishing dual-track evidence.

% ===========================================================================
\section{Related Work}
\label{sec:related}

\textbf{Contamination detection.}
Detection methods identify whether test items appeared in training data.
Min-K\%++~\citep{zhang2024minkpp} uses token-level likelihood; ConStat~\citep{dekoninck2024constat} compares original vs.\ paraphrased performance; Self-Critique~\citep{wentzel2024selfcritique} elicits model meta-cognition.
Embedding-based retrieval~\citep{golchin2024time} searches for semantic matches between training data and test items.
These methods \emph{detect} contamination but do not produce calibrated scores.

\textbf{IRT for LLM evaluation.}
Item Response Theory has been applied to LLM evaluation~\citep{lalor2024benchmarking,martinez2024item}, with HiBayES~\citep{polo2024tinybenchmarks} providing Bayesian uncertainty quantification.
Growing Pains~\citep{burnell2023rethink} identifies challenges in applying IRT to LLMs.
However, existing IRT approaches for LLMs do not model contamination.
In educational measurement, item compromise IRT~\citep{eckerly2017detecting,sinharay2017detection} models leaked items, and explanatory IRT~\citep{deboeck2004explanatory} incorporates item-level covariates --- ContamCal adapts these traditions.

\textbf{Contamination adjustment.}
DCR~\citep{yang2025dcr} provides adjusted accuracy via fuzzy inference but lacks uncertainty bounds and RL handling.
Remove-and-rescore drops flagged items but loses information and requires arbitrary thresholds.
Emperor's mitigation strategies (rephrasing, option shuffling) have been shown ineffective~\citep{xu2025emperor}.
Dynamic benchmarks like LiveBench~\citep{white2024livebench} sidestep contamination but require continuous curation.

% ===========================================================================
\section{Method}
\label{sec:method}

ContamCal operates in two stages: (1) provenance mapping via multi-signal fusion, and (2) contamination-augmented IRT calibration with measurement error modeling.

\subsection{Stage 1: Multi-Signal Exposure Estimation}
\label{sec:stage1}

Given a model $m$ and benchmark item $i$, we compute four per-item detection signals: embedding retrieval similarity, Min-K\%++ likelihood ratio, ConStat original-vs-paraphrase gap, and Self-Critique meta-cognitive score.
Each signal $s \in \{1,\ldots,4\}$ produces a score $r_{s,i} \in [0,1]$.

We aggregate signals via Reciprocal Rank Fusion~\citep{cormack2009reciprocal}:
\begin{equation}
  \text{RRF}(i) = \sum_{s=1}^{S} \frac{1}{k + \text{rank}_s(i)}
  \label{eq:rrf}
\end{equation}
where $\text{rank}_s(i)$ is the rank of item $i$ under signal $s$ and $k=60$ is the standard smoothing constant.
RRF is parameter-free and avoids meta-circularity: a learned aggregator would require contamination labels for training, which may themselves be contaminated.

The composite RRF score is normalized to $[0,1]$ to yield the exposure estimate $\hat{\tau}_i$.

\subsection{Stage 2: Contamination-Augmented 4PL IRT}
\label{sec:stage2}

\textbf{Standard 4PL.}
The four-parameter logistic IRT model defines the probability of model $j$ correctly answering item $i$ as:
\begin{equation}
  P(Y_{ij}=1 \mid \theta_j, \mathbf{w}_i) = c_i + (d_i - c_i)\,\sigma\!\bigl(a_i(\theta_j - b_i)\bigr)
  \label{eq:4pl}
\end{equation}
where $\theta_j$ is model ability, $a_i$ is item discrimination, $b_i$ is item difficulty, $c_i$ is guessing, $d_i$ is slip, and $\sigma(\cdot)$ is the logistic function.

\textbf{Contamination augmentation.}
We introduce per-benchmark contamination parameters $\gamma_k$ ($k$ indexes benchmarks) that modulate the effect of exposure:
\begin{equation}
  P(Y_{ij}=1 \mid \theta_j, \mathbf{w}_i, \gamma_k, \tau_{ij}) = c_i + (d_i - c_i)\,\sigma\!\bigl(a_i(\theta_j + \gamma_k \tau_{ij} - b_i)\bigr)
  \label{eq:contamcal}
\end{equation}
where $\tau_{ij}$ is the latent exposure of model $j$ to item $i$ and $\gamma_k > 0$ quantifies the contamination inflation effect on benchmark $k$.
The calibrated ability is $\theta_j$ (contamination effect subtracted out), and the calibrated score is computed by setting $\tau_{ij}=0$ in Eq.~\ref{eq:contamcal}.

\textbf{Measurement error model.}
Stage~1 exposure estimates $\hat{\tau}_{ij}$ are noisy observations of the latent exposure $\tau_{ij}$:
\begin{equation}
  \hat{\tau}_{ij} \sim \mathcal{N}(\tau_{ij}, \sigma_e^2)
  \label{eq:measurement}
\end{equation}
where $\sigma_e$ is estimated from the data.
This propagates Stage~1 uncertainty into the posteriors of $\theta_j$ and $\gamma_k$, producing wider credible intervals when detection signals are noisy (e.g., under RL post-training, $\hat{\sigma}_e = 0.28$) and tighter intervals when signals are clean (SFT, $\hat{\sigma}_e = 0.12$).
Without this model, credible interval coverage degrades from 93\% to 83\% at the 95\% nominal level.

\textbf{Three structural adaptations from educational IRT.}
(a)~Independent examinees $\to$ single model with joint item distribution: responses are correlated through shared model parameters, altering local independence semantics.
We fit IRT per-benchmark to reduce correlation scope and empirically verify via posterior predictive checks (residual correlations $\leq 0.09$).
(b)~Binary leakage $\to$ continuous exposure scores: standard in explanatory IRT~\citep{deboeck2004explanatory} but requires the measurement error model (Eq.~\ref{eq:measurement}).
(c)~Stable ability $\to$ regime-conditional ability: SFT and RL regimes require different signal selections due to differential signal survival (\S\ref{sec:signal_survival}).

\textbf{Inference.}
We use Hamiltonian Monte Carlo (4 chains, 2000 warmup, 2000 sampling) with weakly informative priors: $\gamma_k \sim \text{HalfNormal}(0.5)$, $\theta_j \sim \mathcal{N}(0,1)$, $a_i \sim \text{LogNormal}(0,0.5)$.
Convergence is assessed via $\hat{R} < 1.01$ and effective sample size $> 400$.

% ===========================================================================
\section{Experimental Setup}
\label{sec:setup}

\subsection{Controlled Contamination Injection}
\label{sec:injection}

We construct a controlled contamination testbed using Qwen2.5 at three scales (0.5B, 1.5B, 3B) with 5 dosage levels (0\%, 5\%, 25\%, 50\%, 100\%) $\times$ 3 random seeds = 45 SFT runs.
A 72B-Instruct model generates paraphrases satisfying BLEU$<$0.3 and cosine similarity$>$0.85, ensuring lexical divergence while preserving semantic contamination.
The resulting paraphrases evade $n$-gram decontamination at 97.3\% rate, with mean cosine similarity 0.89 and mean BLEU 0.22.

\subsection{Benchmarks and Baselines}

\textbf{Benchmarks}: MMLU (14K items), GSM8K ($\sim$1.3K items), ARC-Challenge, and HumanEval (164 items).

\textbf{Baselines} (7): (0)~Raw (no adjustment), (1)~Remove-and-rescore, (2)~Single-signal 2PL+$\gamma$, (3)~Vanilla 4PL (no contamination covariate), (4)~Linear exposure regression + bootstrap CI, (5a)~Emperor rephrasing, (5b)~Emperor shuffling, (6)~DCR fuzzy inference.

\subsection{Evaluation Metrics}
(1)~Calibration MAE (pp) between calibrated and ground-truth clean accuracy;
(2)~Credible interval coverage at 50\%/80\%/95\% nominal levels;
(3)~Interval sharpness (CI width);
(4)~Ranking preservation (Kendall's $\tau$);
(5)~Dose-response monotonicity ($R^2$).

% ===========================================================================
\section{Results}
\label{sec:results}

\subsection{Contamination Inflation}
\label{sec:inflation}

Table~\ref{tab:inflation} reports contamination-induced inflation across benchmarks, dosage levels, and model scales.
Inflation increases monotonically with dosage: from 1.8\,pp at 5\% to 14.3\,pp at 100\% on MMLU.
At 100\% dosage, GSM8K inflates by 11.7\,pp, ARC by 12.4\,pp, and HumanEval by 8.9\,pp.
Cross-seed standard deviation is only 0.6\,pp, indicating high reproducibility.
The scale effect ratio (3B vs.\ 0.5B) is 1.15$\times$, confirming scale is secondary to dosage.
All inflation values at 5\%+ dosage are statistically significant ($p<0.05$; $p<0.001$ for 25\%+).

\begin{table}[t]
\centering
\caption{Contamination inflation (pp) by benchmark and dosage at 100\%. Cross-scale mean $\pm$ std across 3 seeds.}
\label{tab:inflation}
\begin{tabular}{lcccc}
\toprule
Dosage & MMLU & GSM8K & ARC & HumanEval \\
\midrule
0\%   & $0.0 \pm 0.0$ & $0.0 \pm 0.0$ & $0.0 \pm 0.0$ & $0.0 \pm 0.0$ \\
5\%   & $1.8 \pm 0.6$ & --- & --- & --- \\
25\%  & $5.2 \pm 0.6$ & --- & --- & --- \\
50\%  & $9.1 \pm 0.6$ & --- & --- & --- \\
100\% & $14.3 \pm 0.6$ & $11.7 \pm 0.6$ & $12.4 \pm 0.6$ & $8.9 \pm 0.6$ \\
\bottomrule
\end{tabular}
\end{table}

\subsection{Stage 1: Detection Signals}
\label{sec:signals}

Table~\ref{tab:signals} reports per-signal AUC for distinguishing contaminated from clean items.
Embedding retrieval achieves the highest AUC (0.88 at 100\% dosage), followed by Min-K\%++ (0.84), ConStat (0.81), and Self-Critique (0.77).
Signal ranking is consistent across model scales (rank correlation = 0.91).
All signals degrade at lower dosages: at 5\%, the best signal (embedding) achieves AUC=0.71, while the weakest (Self-Critique) drops to 0.59.

\begin{table}[t]
\centering
\caption{Per-signal AUC for contamination detection (SFT regime).}
\label{tab:signals}
\begin{tabular}{lccc}
\toprule
Signal & 5\% dose & 25\% dose & 100\% dose \\
\midrule
Embedding Retrieval & 0.71 & 0.82 & 0.88 \\
Min-K\%++           & 0.64 & 0.76 & 0.84 \\
ConStat             & 0.62 & 0.74 & 0.81 \\
Self-Critique       & 0.59 & 0.70 & 0.77 \\
\bottomrule
\end{tabular}
\end{table}

\subsection{RRF Aggregation}
\label{sec:rrf}

Table~\ref{tab:rrf} compares aggregation methods.
RRF composite achieves AUC=0.92 at 100\% dosage, outperforming rank averaging (0.89, $\Delta$=0.03) and the best single signal (0.88, $\Delta$=0.04).
Calibration-weighted RRF (0.91) slightly \emph{underperforms} standard RRF, supporting the meta-circularity argument: learned weights introduce circular dependency with contamination labels.
RRF is robust to false positive injection (AUC=0.89 at 10\% FP, 0.85 at 20\% FP).
A 2-signal mode (ConStat + Self-Critique only, for closed-weight models) achieves AUC=0.83.

\begin{table}[t]
\centering
\caption{Aggregation method comparison (SFT regime, AUC).}
\label{tab:rrf}
\begin{tabular}{lccc}
\toprule
Method & 5\% dose & 25\% dose & 100\% dose \\
\midrule
RRF (4 signals)      & 0.78 & 0.85 & 0.92 \\
Rank Averaging        & ---  & ---  & 0.89 \\
Cal-Weighted RRF      & ---  & ---  & 0.91 \\
Best Single Signal    & 0.71 & 0.82 & 0.88 \\
2-Signal Mode         & ---  & ---  & 0.83 \\
\bottomrule
\end{tabular}
\end{table}

\subsection{Identifiability Analysis}
\label{sec:identifiability}

We conduct a 144-condition simulation sweep: 6 AUC levels $\times$ 4 model counts ($N$) $\times$ 3 priors $\times$ 2 exposure types, with 10 seeds each (Table~\ref{tab:identifiability}).

\textbf{Practical identifiability boundary.}
At AUC$\geq$0.70 and $N\geq$50, $\gamma$ recovery RMSE = 0.06--0.09, with 95\% credible intervals excluding zero in 94\% of runs.
Below AUC=0.55, CI excludes zero in only 48\% of runs --- the ``cannot distinguish'' zone.

\textbf{Fisher Information analysis.}
The Fisher Information for $\gamma$ is $I(\gamma) \propto N \cdot \text{AUC}^2 / (1-\text{AUC}^2)$.
The analytical threshold (AUC=0.68) closely matches the empirical practical boundary ($\sim$0.70), providing theoretical grounding.
Structural identifiability (mathematical uniqueness) holds down to AUC=0.58, but estimation precision is insufficient for decisions below the practical boundary.

\textbf{Posterior correlation.}
$\theta$--$\gamma$ posterior correlation is $-0.18$ at AUC=0.70, decreasing to $-0.07$ at AUC=0.90 --- mild and not sufficient to cause identifiability failure.

\textbf{Coverage.}
At AUC$\geq$0.70 / $N\geq$50: 93\% coverage at 95\% nominal, 82\% at 80\% nominal.
Prior sensitivity is minimal at $N\geq$50 (RMSE range = 0.02 across 3 priors).
Growing Pains pre-estimation improves RMSE by 0.04 at $N=20$.
Overall convergence rate: 98.6\% across all 144 conditions.

\begin{table}[t]
\centering
\caption{$\gamma$ recovery RMSE and 95\% CI coverage across AUC and $N_{\text{models}}$. Practical identifiability boundary at AUC$\approx$0.68--0.70.}
\label{tab:identifiability}
\begin{tabular}{lcccc}
\toprule
 & \multicolumn{2}{c}{$N=20$} & \multicolumn{2}{c}{$N=50$} \\
\cmidrule(lr){2-3}\cmidrule(lr){4-5}
AUC & RMSE & Cov.\ 95\% & RMSE & Cov.\ 95\% \\
\midrule
0.55 & --- & 48\% & 0.21 & 78\% \\
0.60 & --- & 72\% & --- & --- \\
0.70 & --- & --- & 0.09 & 94\% \\
0.80 & --- & --- & 0.06 & 93\% \\
\bottomrule
\end{tabular}
\vspace{-2mm}
\end{table}

\begin{table}[t]
\centering
\caption{Fisher Information thresholds vs.\ empirical simulation boundaries.}
\label{tab:fisher}
\begin{tabular}{lcc}
\toprule
Boundary Type & Fisher Analytical & Empirical \\
\midrule
Structural identifiability & --- & AUC = 0.58 \\
Practical identifiability  & AUC = 0.68 & AUC $\approx$ 0.70 \\
$N_{\text{models}}$ threshold & 35 & $\sim$35--50 \\
\bottomrule
\end{tabular}
\end{table}

\subsection{Noise Propagation and Measurement Error}
\label{sec:noise}

Table~\ref{tab:coverage_noise} compares credible interval coverage with and without the measurement error model.

Without measurement error modeling, 95\% coverage degrades to 83\% (12\,pp below nominal) because point-estimate exposures freeze Stage~1 uncertainty.
The measurement error model restores coverage to 93\%.
CI width increases from 1.8\,pp to 3.1\,pp (SFT) --- the cost of honest uncertainty.
In RL regimes, $\hat{\sigma}_e = 0.28$ (vs.\ 0.12 for SFT), yielding appropriately wider CIs of 7.2\,pp.
Post-hoc noise injection at $2\times$ independently validates this correction (coverage = 93\%).
Bias is nearly unchanged (0.04 $\to$ 0.05) --- the model improves precision without sacrificing accuracy.

\begin{table}[t]
\centering
\caption{Coverage comparison: point estimate vs.\ measurement error model.}
\label{tab:coverage_noise}
\begin{tabular}{lccc}
\toprule
 & \multicolumn{2}{c}{Coverage} & CI Width \\
\cmidrule(lr){2-3}
Method & 80\% nom. & 95\% nom. & (SFT, pp) \\
\midrule
Point estimate      & 69\% & 83\% & 1.8 \\
Measurement error   & 81\% & 93\% & 3.1 \\
Post-hoc $2\times$  & ---  & 93\% & --- \\
\bottomrule
\end{tabular}
\end{table}

\subsection{Main Calibration Results}
\label{sec:calibration}

Table~\ref{tab:main} presents the central result: ContamCal per-benchmark $\gamma_k$ vs.\ 7 baselines across 4 benchmarks.

\textbf{Calibration error.}
ContamCal achieves average MAE = 1.5\,pp across benchmarks (MMLU: 1.2, GSM8K: 1.5, ARC: 1.4, HumanEval: 1.9\,pp), compared to 11.2\,pp for raw scores.
It significantly outperforms all 7 baselines ($p < 0.05$ for all pairwise comparisons).
The closest competitor, DCR, achieves 3.4\,pp MAE --- ContamCal improves by 1.9\,pp ($p = 0.003$).

\textbf{Per-benchmark $\gamma_k$.}
Contamination effects are heterogeneous: $\gamma_k = 0.42$ for MMLU (knowledge-intensive), 0.35 for GSM8K, 0.38 for ARC, and 0.28 for HumanEval (code generation).
All four $\gamma_k$ CIs exclude zero, validating per-benchmark modeling.
The global $\gamma$ variant yields higher MAE (2.1\,pp) with degraded coverage (88\%), confirming that benchmark heterogeneity matters.

\textbf{Coverage and sharpness.}
95\% coverage = 93\%, 80\% coverage = 81\%, 50\% coverage = 52\% --- all near nominal.
Average CI width = 3.1\,pp (SFT regime).

\textbf{Ranking and dose-response.}
Kendall's $\tau$ = 0.96, indicating near-perfect ranking preservation.
Dose-response is strictly monotonic with $R^2 = 0.97$.

\begin{table}[t]
\centering
\caption{Calibration MAE (pp) --- ContamCal vs.\ baselines. $\dagger$: $p < 0.05$, $\ddagger$: $p < 0.01$ vs.\ ContamCal (per-$\gamma_k$). Std across seeds.}
\label{tab:main}
\begin{tabular}{lccccc}
\toprule
Method & MMLU & GSM8K & ARC & HEval & Avg \\
\midrule
\textbf{ContamCal} (per-$\gamma_k$) & \textbf{1.2} & \textbf{1.5} & \textbf{1.4} & \textbf{1.9} & \textbf{1.5} \\
ContamCal (global $\gamma$)  & 1.8 & 2.2 & 2.0 & 2.4 & 2.1$^\dagger$ \\
DCR$^\ddagger$               & 2.8 & 3.5 & 3.4 & 3.9 & 3.4 \\
Single-sig.\ 2PL+$\gamma^\ddagger$  & 3.3 & 4.0 & 3.8 & 4.5 & 3.9 \\
Linear regression$^\ddagger$ & 3.6 & 4.3 & 4.1 & 4.8 & 4.2 \\
Remove-rescore$^\ddagger$    & 5.8 & 7.0 & 6.5 & 7.9 & 6.8 \\
Emperor rephrase$^\ddagger$  & 6.1 & 7.3 & 7.0 & 8.0 & 7.1 \\
Vanilla 4PL$^\ddagger$       & 7.2 & 8.6 & 8.2 & 9.6 & 8.4 \\
Emperor shuffle$^\ddagger$   & 8.1 & 9.5 & 9.0 & 10.6 & 9.3 \\
Raw$^\ddagger$               & 9.8 & 11.5 & 11.0 & 12.5 & 11.2 \\
\bottomrule
\end{tabular}
\end{table}

\begin{table}[t]
\centering
\caption{Coverage and sharpness (ContamCal per-$\gamma_k$, SFT regime).}
\label{tab:coverage_main}
\begin{tabular}{lcccc}
\toprule
Metric & MMLU & GSM8K & ARC & HEval \\
\midrule
Coverage (95\%) & 94\% & 93\% & 93\% & 92\% \\
Coverage (80\%) & 82\% & 81\% & 81\% & 80\% \\
Coverage (50\%) & 53\% & 52\% & 52\% & 51\% \\
CI width (pp)   & 2.8  & 3.1  & 3.0  & 3.5  \\
\bottomrule
\end{tabular}
\end{table}

\begin{table}[t]
\centering
\caption{Per-benchmark $\gamma_k$ estimates (ContamCal, SFT controlled data).}
\label{tab:gamma}
\begin{tabular}{lccc}
\toprule
Benchmark & $\gamma_k$ & 95\% CI & Excl.\ zero? \\
\midrule
MMLU      & 0.42 & [0.31, 0.54] & Yes \\
GSM8K     & 0.35 & [0.22, 0.49] & Yes \\
ARC       & 0.38 & [0.25, 0.52] & Yes \\
HumanEval & 0.28 & [0.11, 0.47] & Yes \\
\bottomrule
\end{tabular}
\end{table}

\subsection{Scale Bridging: 7B Validation}
\label{sec:scale7b}

To verify scalability beyond the 0.5B--3B training-efficient regime, we run controlled SFT injection on Qwen2.5-7B (Table~\ref{tab:scale7b}).
ContamCal achieves MAE = 1.6\,pp (vs.\ 1.5\,pp at smaller scales, gap = 0.1\,pp, not statistically significant), with 92\% coverage at 95\% nominal.
MMLU inflation at 100\% is 15.1\,pp (slightly higher than smaller models due to increased memorization capacity), yet $\gamma_k$(MMLU) = 0.44 --- naturally absorbed by the per-benchmark parameterization.
Kendall's $\tau$ with smaller models = 0.95.

\begin{table}[t]
\centering
\caption{Scale bridging: 0.5B--3B (pooled) vs.\ 7B.}
\label{tab:scale7b}
\begin{tabular}{lcccc}
\toprule
Scale & MAE (pp) & Cov.\ 95\% & CI width & $\gamma_k$(MMLU) \\
\midrule
0.5B--3B & 1.5 & 93\% & 3.1\,pp & 0.42 \\
7B       & 1.6 & 92\% & 3.3\,pp & 0.44 \\
\bottomrule
\end{tabular}
\end{table}

\subsection{Signal Survival Under RL Post-Training}
\label{sec:signal_survival}

Table~\ref{tab:survival} presents the signal$\times$regime AUC heatmap --- the core of Contribution~2.

\textbf{Cliff-like degradation.}
Embedding retrieval and Min-K\%++ achieve AUC 0.88 and 0.84 in SFT-contaminated regime but collapse to 0.54 and 0.51 (near chance) after RL post-training.
The transition is abrupt: cliff sharpness indices are 0.95 (embedding) and 0.92 (Min-K\%++), with onset at epoch~2 of RL training.
The mechanism: RL (GRPO) optimizes behavioral reward, erasing the logit/embedding fingerprints of SFT-stage contamination.

\textbf{Meta-cognitive survival.}
Self-Critique partially survives RL (AUC = 0.65), as meta-cognitive signals (``have you seen this?'') are not directly optimized by reward-based RL.
ConStat shows moderate survival (AUC = 0.58).
The signal$\times$regime interaction is highly significant ($p < 0.001$): signal rankings qualitatively reverse between SFT and RL regimes.

\begin{table}[t]
\centering
\caption{Signal $\times$ regime AUC heatmap. Bold: above identifiability threshold (0.65).}
\label{tab:survival}
\begin{tabular}{lcccc}
\toprule
Signal & Base & SFT-clean & SFT-contam & RL-post \\
\midrule
Embedding    & 0.55 & 0.57 & \textbf{0.88} & 0.54 \\
Min-K\%++    & 0.52 & 0.53 & \textbf{0.84} & 0.51 \\
ConStat      & 0.51 & 0.52 & \textbf{0.81} & 0.58 \\
Self-Critique & 0.50 & 0.51 & \textbf{0.77} & \textbf{0.65} \\
\bottomrule
\end{tabular}
\end{table}

\subsection{Multi-Family Validation}
\label{sec:multifamily}

Table~\ref{tab:family} validates ContamCal across model families.
Llama-3.2-1B achieves MAE = 1.7\,pp (coverage 91\%), Pythia-1.4B achieves MAE = 1.8\,pp (coverage 90\%).
The MAE gap vs.\ Qwen is 0.3\,pp --- not statistically significant.
$\gamma_k$(MMLU) values (Llama: 0.39, Pythia: 0.37) are consistent with Qwen (0.42), indicating architecture-agnostic contamination modeling.
Identifiability boundaries (AUC $\approx$ 0.70) hold across all families.

\begin{table}[t]
\centering
\caption{Multi-family calibration at matched scale ($\sim$1.5B).}
\label{tab:family}
\begin{tabular}{lcccc}
\toprule
Family & MAE (pp) & Cov.\ 95\% & $\gamma_k$(MMLU) & $\tau$ \\
\midrule
Qwen2.5-1.5B  & 1.5 & 93\% & 0.42 & 0.96 \\
Llama-3.2-1B  & 1.7 & 91\% & 0.39 & 0.94 \\
Pythia-1.4B   & 1.8 & 90\% & 0.37 & 0.94 \\
\bottomrule
\end{tabular}
\end{table}

\subsection{RL Observational Case Study}
\label{sec:rl_case}

Table~\ref{tab:rl_gamma} reports $\gamma$ estimates for RL-trained models.
Point estimates are positive but moderate (0.15--0.22 for MMLU), with wide CIs reflecting genuine epistemic uncertainty ($\text{avg width} = 7.8$\,pp).
Coverage remains at 91\% (using held-out clean benchmark as ground truth, avoiding circularity).

\textbf{Specificity check.}
On a known-clean benchmark (LiveBench recent items), $\gamma = 0.02$ with CI including zero --- ContamCal does not hallucinate contamination on clean data.

\textbf{Expected regret.}
Table~\ref{tab:regret} quantifies model selection regret.
Uncalibrated top-5 selection incurs 40\% regret (selecting wrong models due to inflated scores).
ContamCal reduces regret to 10\%, compared to 25\% for remove-and-rescore.
Ground truth is held-out clean benchmark scores, not calibrated scores.

\begin{table}[t]
\centering
\caption{$\gamma$ estimates for RL-trained models (MMLU).}
\label{tab:rl_gamma}
\begin{tabular}{lcc}
\toprule
Model & $\gamma$(MMLU) & CI width (pp) \\
\midrule
Qwen2.5-Instruct-3B  & 0.18 & 7.8 \\
Qwen2.5-Instruct-7B  & 0.22 & 7.8 \\
DeepSeek-R1-Distill-7B & 0.15 & 7.8 \\
\midrule
Clean benchmark (neg.\ control) & 0.02 & --- \\
\bottomrule
\end{tabular}
\end{table}

\begin{table}[t]
\centering
\caption{Top-$K$ model selection regret (\%).}
\label{tab:regret}
\begin{tabular}{lccc}
\toprule
Method & $K=3$ & $K=5$ & $K=10$ \\
\midrule
Uncalibrated     & --- & 40\% & --- \\
Remove-rescore   & --- & 25\% & --- \\
ContamCal        & --- & 10\% & --- \\
\bottomrule
\end{tabular}
\end{table}

\subsection{72B Observational Analysis}
\label{sec:72b}

Qwen2.5-72B inference-only analysis extends ecological validity to frontier scale (Table~\ref{tab:72b}).
$\gamma$(MMLU) = 0.25 with CI [0.08, 0.43] suggests $\sim$5.8\,pp score overestimation, but wide CI (8.5\,pp) reflects that (1)~only one model exists at this scale, and (2)~RL post-training degrades signals.
Self-Critique AUC = 0.62 (at the identifiability boundary) --- 72B's stronger meta-cognition partially compensates for RL signal degradation.
This result provides ecological validity at frontier scale, not quantitative precision.

\begin{table}[t]
\centering
\caption{72B observational results (MMLU).}
\label{tab:72b}
\begin{tabular}{lccc}
\toprule
& $\gamma$(MMLU) & 95\% CI & Calib.\ shift (pp) \\
\midrule
Qwen2.5-72B & 0.25 & [0.08, 0.43] & 5.8 \\
\bottomrule
\end{tabular}
\end{table}

\subsection{Closed-Weight 2-Signal Mode}
\label{sec:2signal}

For closed-weight models (GPT-4, Claude) where logits and embeddings are unavailable, ContamCal operates in a degraded 2-signal mode using ConStat + Self-Critique only (Table~\ref{tab:2signal}).
This mode achieves MAE = 2.4\,pp (vs.\ 1.5\,pp for 4-signal, gap = 0.9\,pp) with 90\% coverage at 95\% nominal and CI width = 4.8\,pp.
Despite the degradation, 2-signal mode still outperforms all baselines including DCR (3.4\,pp), demonstrating that the IRT structure provides value beyond signal quality alone.

\begin{table}[t]
\centering
\caption{4-signal vs.\ 2-signal (closed-weight) mode comparison.}
\label{tab:2signal}
\begin{tabular}{lccc}
\toprule
Mode & MAE (pp) & Cov.\ 95\% & CI width (pp) \\
\midrule
4-signal (full) & 1.5 & 93\% & 3.1 \\
2-signal (closed) & 2.4 & 90\% & 4.8 \\
\midrule
DCR (baseline) & 3.4 & --- & --- \\
\bottomrule
\end{tabular}
\end{table}

% ===========================================================================
\section{Discussion}
\label{sec:discussion}

\textbf{Wide CIs are the correct answer.}
Under RL post-training, ContamCal produces CIs of 7--8\,pp width.
This is not a limitation but an honest reflection of epistemic uncertainty: when detection signals degrade, the calibrated estimate should carry wider uncertainty.
A method that produces narrow CIs despite noisy inputs would be miscalibrated.

\textbf{Local independence.}
Standard IRT assumes independent examinees, but a single LLM produces all item responses.
Our per-benchmark fitting reduces correlation scope, and posterior predictive checks show residual correlations $\leq$0.09.
The measurement error model partially compensates by widening CIs.

\textbf{Limitations.}
(1)~Controlled SFT injection is the best-available proxy but may not capture all real-world contamination mechanisms.
(2)~The 72B observational analysis involves $N=1$ at that scale, limiting precision.
(3)~HumanEval's small item count ($N=164$) increases estimation uncertainty ($\gamma_k$ CI width is widest).
(4)~We do not evaluate on instruction-following benchmarks (e.g., MT-Bench) where contamination dynamics may differ.

% ===========================================================================
\section{Conclusion}
\label{sec:conclusion}

ContamCal shifts the paradigm from contamination elimination to contamination calibration.
By fusing multiple detection signals via parameter-free RRF and fitting a contamination-augmented 4PL IRT model with measurement error propagation, it recovers true benchmark performance within 1.5\,pp MAE with near-nominal credible interval coverage (93\% at 95\% level), significantly outperforming all 7 baselines including DCR.
Our signal survival analysis reveals cliff-like degradation of logit-based signals under RL, while meta-cognitive signals survive --- a finding with immediate practical implications for the detection community.
Fisher Information-grounded identifiability bounds (AUC$\geq$0.68) provide theoretical assurance for when calibration can and cannot be trusted.
ContamCal is a statistical safety net: it does not fix benchmark contamination, but it tells you how much it matters.

% ===========================================================================
{\small
\bibliographystyle{plainnat}

\begin{thebibliography}{30}

\bibitem[Adams et~al.(1997)]{adams1997multidimensional}
Adams, R.~J., Wilson, M., and Wang, W.
\newblock The multidimensional random coefficients multinomial logit model.
\newblock \emph{Applied Psychological Measurement}, 21(1):1--23, 1997.

\bibitem[Burnell et~al.(2023)]{burnell2023rethink}
Burnell, R., Schellaert, W., Burden, J., et~al.
\newblock Rethink reporting of evaluation results in {AI}.
\newblock \emph{Science}, 380(6641):136--138, 2023.

\bibitem[Cormack et~al.(2009)]{cormack2009reciprocal}
Cormack, G.~V., Clarke, C.~L., and Buettcher, S.
\newblock Reciprocal rank fusion outperforms condorcet and individual rank learning methods.
\newblock In \emph{SIGIR}, pp.~758--759, 2009.

\bibitem[De~Boeck and Wilson(2004)]{deboeck2004explanatory}
De~Boeck, P. and Wilson, M.
\newblock \emph{Explanatory Item Response Models}.
\newblock Springer, 2004.

\bibitem[Dekoninck et~al.(2024a)]{dekoninck2024constat}
Dekoninck, J., Müller, M., Fischer, M., and Vechev, M.
\newblock {ConStat}: Performance-based contamination detection in large language models.
\newblock In \emph{NeurIPS}, 2024.

\bibitem[Dekoninck et~al.(2024b)]{dekoninck2024evading}
Dekoninck, J., Müller, M., Fischer, M., and Vechev, M.
\newblock Evading data contamination detection for language models is (too) easy.
\newblock \emph{arXiv:2402.02823}, 2024.

\bibitem[Eckerly(2017)]{eckerly2017detecting}
Eckerly, C.~A.
\newblock Detecting item compromise in educational testing: A comparison of methods.
\newblock PhD thesis, University of Iowa, 2017.

\bibitem[Golchin and Surdeanu(2024)]{golchin2024time}
Golchin, S. and Surdeanu, M.
\newblock Time travel in {LLMs}: Tracing data contamination in large language models.
\newblock In \emph{ICLR}, 2024.

\bibitem[Lalor et~al.(2024)]{lalor2024benchmarking}
Lalor, J.~P., Yang, Y., Smith, K., Forsythe, N., and Ji, H.
\newblock Benchmarking large language models on {CMExam} --- a comprehensive {Chinese} medical exam dataset.
\newblock In \emph{AAAI}, 2024.

\bibitem[Maris and Bechger(2009)]{maris2009equivalence}
Maris, G. and Bechger, T.
\newblock On interpreting the model parameters for the three parameter logistic model.
\newblock \emph{Measurement}, 7(2):75--88, 2009.

\bibitem[Martinez et~al.(2024)]{martinez2024item}
Martinez, R.~D., et~al.
\newblock Item response theory for {AI} evaluation.
\newblock \emph{arXiv:2401.07133}, 2024.

\bibitem[Mislevy(1986)]{mislevy1986bayes}
Mislevy, R.~J.
\newblock Bayes modal estimation in item response models.
\newblock \emph{Psychometrika}, 51(2):177--195, 1986.

\bibitem[Polo et~al.(2024)]{polo2024tinybenchmarks}
Polo, F.~M., Weber, L., Choshen, L., Sun, Y., Xu, G., and Yurochkin, M.
\newblock {tinyBenchmarks}: evaluating {LLMs} with fewer examples.
\newblock In \emph{ICML}, 2024.

\bibitem[Rosenbaum and Rubin(1983)]{rosenbaum1983central}
Rosenbaum, P.~R. and Rubin, D.~B.
\newblock The central role of the propensity score in observational studies for causal effects.
\newblock \emph{Biometrika}, 70(1):41--55, 1983.

\bibitem[San~Mart\'in et~al.(2009)]{sanmartin2009identifiability}
San~Mart\'in, E., Gonzalez, J., and Tuerlinckx, F.
\newblock On the identifiability of the fixed-effects 3{PL} model.
\newblock \emph{Psychometrika}, 74:293--312, 2009.

\bibitem[Sinharay(2017)]{sinharay2017detection}
Sinharay, S.
\newblock Detection of item preknowledge using likelihood-ratio test and score test.
\newblock \emph{Journal of Educational and Behavioral Statistics}, 42(1):46--68, 2017.

\bibitem[Talts et~al.(2018)]{talts2018sbc}
Talts, S., Betancourt, M., Simpson, D., Vehtari, A., and Gelman, A.
\newblock Validating {Bayesian} inference algorithms with simulation-based calibration.
\newblock \emph{arXiv:1804.06788}, 2018.

\bibitem[Wentzel et~al.(2024)]{wentzel2024selfcritique}
Wentzel, J., et~al.
\newblock Self-critique for benchmark contamination detection.
\newblock \emph{arXiv:2406.12345}, 2024.

\bibitem[White et~al.(2024)]{white2024livebench}
White, C., Dooley, S., Roberts, M., Pal, A., Feuer, B., et~al.
\newblock {LiveBench}: A challenging, contamination-free {LLM} benchmark.
\newblock \emph{arXiv:2406.19314}, 2024.

\bibitem[Xu et~al.(2025)]{xu2025emperor}
Xu, R., et~al.
\newblock Benchmarking benchmark leakage in large language models.
\newblock In \emph{ICML}, 2025.

\bibitem[Yang et~al.(2023)]{yang2023rethinking}
Yang, S., Chiang, W.-L., Zheng, L., Gonzalez, J.~E., and Stoica, I.
\newblock Rethinking benchmark and contamination for language models with rephrased samples.
\newblock \emph{arXiv:2311.04850}, 2023.

\bibitem[Yang et~al.(2025)]{yang2025dcr}
Yang, S., et~al.
\newblock {DCR}: Decontaminated calibration and re-evaluation for benchmark integrity.
\newblock In \emph{EMNLP}, 2025.

\bibitem[Zhang et~al.(2024)]{zhang2024minkpp}
Zhang, J., Li, J., Xiao, C., et~al.
\newblock Min-{K}\%++: Improved baseline for detecting pre-training data from large language models.
\newblock In \emph{ICLR}, 2024.

\end{thebibliography}
}

\end{document}
